%% =====================================================================
%% WS-4 Annual Review Paper #1 — npj Digital Medicine submission
%%
%% Target venue: npj Digital Medicine (Nature Portfolio; clinical-AI scope)
%% Fallback:     Medical Image Analysis (review article track)
%% Template:     Nature-family. We use article in development; switch to
%%               the Nature Portfolio sn-jnl.cls at submission.
%%
%% Status: v0.1 working draft — PROTOCOL ONLY. Year-1 results are
%%         placeholders since the competition itself does not exist
%%         until D9+730. This draft establishes the manuscript template
%%         that will be cloned every year thereafter as the recurring
%%         "State of the Field" annual review.
%%
%% Build:  pdflatex manuscript && bibtex manuscript && pdflatex manuscript
%% =====================================================================

\documentclass[twoside,11pt]{article}

\usepackage[utf8]{inputenc}
\usepackage[T1]{fontenc}
\usepackage{lmodern}
\usepackage{microtype}
\usepackage[margin=1in]{geometry}
\usepackage{graphicx}
\usepackage{booktabs}
\usepackage{multirow}
\usepackage{amsmath,amssymb}
\usepackage{xcolor}
\usepackage[colorlinks=true,linkcolor=blue,citecolor=blue,urlcolor=blue]{hyperref}
\usepackage{xurl}
\usepackage[round]{natbib}
\usepackage{authblk}

\newcommand{\todo}[1]{\textcolor{red}{\textbf{[TODO: #1]}}}
\newcommand{\todoTable}[2]{%
  \begin{table}[ht]\centering%
    \fbox{\parbox{0.85\linewidth}{\centering\todo{TABLE: #2}}}%
    \caption{#2 (placeholder; populate when data lands).}\label{#1}%
  \end{table}}
\newcommand{\todoFigure}[2]{%
  \begin{figure}[ht]\centering%
    \fbox{\parbox{0.85\linewidth}{\centering\todo{FIGURE: #2}}}%
    \caption{#2 (placeholder; populate when data lands).}\label{#1}%
  \end{figure}}

\title{\textbf{State of the Field: Year 1 of the PWM Low-Dose CT Challenge}}

\author[1,2]{\textit{Authors to be confirmed at submission}}
\affil[1]{University of Texas Southwestern Medical Center, Dallas, TX, USA}
\affil[2]{PWM Protocol Foundation}

\date{Working draft v0.1 --- \today \\ {\small (will be re-cut annually as new yearly editions; this draft is the Year-1 template.)}}

\begin{document}
\maketitle

%% =====================================================================
\begin{abstract}
\textbf{Background.} Low-dose CT reconstruction has accumulated years of method papers without producing a citation standard against which longitudinal progress can be measured. A practitioner deciding which method to deploy on a new vendor's data, or a radiologist asking which dose reduction is now diagnostically safe for a specific clinical task, has no shared reference.

\textbf{Methods.} The PWM Low-Dose CT Challenge is a permanent, content-addressed, citation-standard competition hosted on the PWM blockchain registry, scored against the multi-vendor PWM-LDCT benchmark~\citep{ws1dataset} using the signal-equivalence framework of~\citep{ws2framework}. Submissions take the form of sandboxed Docker RunBundles whose published results are re-executed and bit-identically verified by the scoring service. We review the first year of the challenge: \todo{N submissions} community entries from \todo{N institutions} institutions in \todo{N countries} countries, including \todo{N} industrial vendor entries.

\textbf{Findings.} We characterize Year-1 submissions across five axes: reconstruction quality (PSNR / SSIM / LPIPS), downstream task performance (lung-nodule detection AUC), cross-vendor generalization, signal-equivalence credential coverage at multiple dose ratios, and open-source license posture. Three findings: (i) the field has converged on \todo{trend}; (ii) credential coverage at $r = 0.25$ reached \todo{X\%} but $r = 0.10$ remains an open problem; (iii) cross-vendor generalization remains the largest unsolved gap, median PSNR drop \todo{X dB}.

\textbf{Significance for clinical practice.} For the first time, a radiologist or physicist can consult Table~\ref{tab:clinical_impact_summary} to look up the most aggressive open-source-credentialed dose reduction for a specific task and subpopulation. This paper inaugurates an annual review series~\citep{ws4yearN_template} that will recur each year, citing every submission from that year and providing a longitudinal snapshot of low-dose CT reconstruction as the field evolves.
\end{abstract}

%% =====================================================================
\section{Introduction}

The deep-learning era of low-dose CT reconstruction has produced a steady stream of incremental method papers without producing a citation standard. A new method published in 2023 typically reports a PSNR improvement of fractions of a decibel over the previous best, on the same dataset (the Mayo Clinic AAPM 2016 cohort~\citep{aapm2016challenge}), using the same evaluation metrics, with little or no cross-vendor evaluation, no per-pixel uncertainty quantification, and a license that varies from Apache 2.0 to ``available on request.'' A practitioner picking a method to deploy on data from a different scanner has no clear choice; a regulator evaluating an FDA submission has no shared reference; a researcher writing a follow-on paper cannot tell whether the previous-best number is from a comparable evaluation protocol.

The PWM Low-Dose CT Challenge addresses these structural gaps by combining three artifacts that are individually well-established but rarely co-located: a \textbf{multi-vendor, paired-dose, on-chain-anchored benchmark dataset}~\citep{ws1dataset}; a \textbf{probabilistic, task-conditional, modality-general signal-equivalence framework}~\citep{ws2framework}; and a \textbf{permanent competition} whose submissions, evaluation, and scoring are reproducible end-to-end via PWM RunBundles~\citep{pwmregistry}. Submissions take the form of a sandboxed Docker container with a pinned eval entry point and a published \texttt{results.json}; the challenge's scoring service re-executes the container, verifies the published results against the live execution, and posts the verified score to a public leaderboard. The L3 dataset spec and the L2 framework spec are content-addressed on the PWM blockchain, so a result tied to \texttt{pwm\_l3\_spec:pwm\_low\_dose\_ct\_benchmark:1.0.0} and \texttt{pwm\_l2\_spec:dose\_equivalence\_v1:0.1.0} is bit-identically reproducible across years even as the dataset is later extended.

\paragraph{Why now.} Three concurrent developments make a content-addressed annual review newly possible. First, multi-vendor paired-dose CT data is now publicly available at scale through the PWM-LDCT benchmark~\citep{ws1dataset}, ending the decade-long monopoly of single-vendor (Siemens) benchmarks. Second, regulatory bodies including the FDA~\citep{fda_aiml_2024} are explicitly asking for subgroup-stratified, version-pinned performance evidence for AI/ML-based medical software; the signal-equivalence framework~\citep{ws2framework} provides the credential format that operationalizes this requirement. Third, public blockchain registries with sub-second verification of cryptographic hashes are now stable infrastructure, making content-addressed benchmark anchoring practical for the first time. The PWM challenge combines all three.

This paper is the inaugural annual review of the challenge. It is the first of an intended sequence: each year, an analogous paper will summarize the cohort of new submissions, characterize cross-method trends, perform clinical-impact analysis (which dose levels are now diagnostically equivalent for which subpopulations), and identify the open problems that should define the next year's work. The recurring annual structure --- with consistent benchmark, framework, and reporting standards across years --- is the mechanism by which the field can finally measure longitudinal progress: by Year 5 of the series (2032), the cumulative leaderboard will contain $\sim$\todo{N$_{\textrm{5-yr}}$} cited methods, each evaluable against bit-identical benchmark and framework versions, producing a longitudinal performance trace no single static review can construct.

\paragraph{What makes this review series distinctive.} Reviews of low-dose CT reconstruction have appeared periodically~\citep{wang2020deeplearningreview}, but they share three limitations that the present series is structurally designed to avoid:

\begin{enumerate}
    \item \textbf{Static snapshots.} Prior reviews are written once. A method published two years later is not visible until the next ad-hoc review. The PWM annual cadence guarantees a yearly cohort summary.
    \item \textbf{Anchorless evaluation.} Prior reviews compare numbers reported on differently-cleaned versions of the same dataset, with no mechanism to verify that those numbers were produced under comparable evaluation protocols. The PWM challenge ties every cited number to a content-addressed benchmark hash and a content-addressed framework hash, so cross-year comparisons are bit-identically reproducible.
    \item \textbf{Pixel-wise metrics only.} Prior reviews emphasize PSNR / SSIM. The PWM challenge requires every submission to report a 5-tuple credential under the signal-equivalence framework~\citep{ws2framework}, which ties the comparison to a specified clinical task, subpopulation, margin, and confidence level.
\end{enumerate}

\paragraph{Contributions of this Year-1 review.}
\begin{enumerate}
    \item A summary of the year-1 submission cohort across methodology classes, institutions, and reported metrics (Section~\ref{sec:year1_cohort}).
    \item Cross-method trends across reconstruction quality, downstream task performance, cross-vendor generalization, and uncertainty quantification (Section~\ref{sec:trends}).
    \item Clinical-impact analysis: which dose levels are now signal-equivalent under the framework of~\citep{ws2framework} for which clinical tasks (Section~\ref{sec:clinical_impact}).
    \item A named list of open problems we believe define Year 2 (Section~\ref{sec:open_problems}).
    \item A protocol for the year-2 challenge cycle (Section~\ref{sec:year2}).
\end{enumerate}

%% =====================================================================
\section{Review methods}
\label{sec:review_methods}

\subsection{Scope of this review}

This review covers all submissions to the PWM Low-Dose CT Challenge between the public launch on \todo{launch date} and the Year-1 close on \todo{close date}. ``Submission'' means a PWM RunBundle that (a) was submitted to the leaderboard scoring service, (b) declared at least one 5-tuple signal-equivalence credential per~\citep{ws2framework}, and (c) was accompanied by a permissively-licensed or research-licensed code release. Submissions retracted by their authors before the Year-1 close were excluded; the retraction itself is recorded in the leaderboard's public history but is not analyzed in this review.

\subsection{Inclusion and exclusion criteria}

Submissions \emph{included} in the cross-method analysis:
\begin{itemize}
    \item Reported at least the canonical 5-tuple credential at $r = 0.25$ for chest lung-nodule detection.
    \item Passed bit-identical re-execution verification by the scoring service within at most three resubmission cycles.
    \item Released code and weights at a stable URL (Apache 2.0, MIT, BSD, or equivalent permissive license; or a documented research-only license with reproducibility commitment).
\end{itemize}

Submissions \emph{excluded} from cross-method analysis but recorded in the leaderboard public log:
\begin{itemize}
    \item Failed verification after three resubmission cycles (\todo{N} submissions).
    \item Reported credentials only at non-canonical $(r, T)$ combinations not analyzed here (\todo{N} submissions; their published credentials remain valid).
    \item Submitted after the Year-1 cutoff (\todo{N} submissions; these will be analyzed in the Year-2 review).
\end{itemize}

The full submission-flow figure (Figure~\ref{fig:review_flow}) summarizes inclusion and exclusion at each step.

\todoFigure{fig:review_flow}{PRISMA-style submission-flow diagram. The rendered figure shows: (top) all RunBundles submitted to the scoring service during the Year-1 window with count; (middle) verification pass / resubmission / unrecoverable-failure outcomes; (bottom) inclusion in cross-method analysis vs leaderboard-log-only. Numbers populated from the scoring service's audit log at Year-1 close.}

\subsection{Conflict of interest in challenge organization}

This review is co-authored by members of the team that organizes the PWM Low-Dose CT Challenge. We disclose the following structural conflicts and the safeguards adopted:

\begin{itemize}
    \item \textbf{Organizer submissions.} The PWM reference reconstruction method~\citep{ws3reference} is the seed entry in the leaderboard. To avoid self-citation bias, the analyses in this review present the reference method's credentials transparently alongside community submissions but the leaderboard headline ranking is based on the canonical task metric (lung-nodule AUC at $r = 0.25$), which the reference method does \emph{not} top.
    \item \textbf{Authorship of submitted methods.} Co-authors of this review who are also authors of a submitted method are explicitly listed in the Acknowledgments section by initials, and their submissions are flagged in the leaderboard public log with \texttt{author\_of\_challenge\_review: true}.
    \item \textbf{Reviewer assignment for cited submissions.} This review summarizes but does not peer-review individual submissions; submission-level peer review is the responsibility of the submission's own publication venue.
\end{itemize}

\subsection{Reproducibility of this review}

The numerical content of this review (counts in Tables and Figures) is generated by a single script (\texttt{review/year1\_analysis.py}) that ingests the scoring service's audit log and produces every cited number deterministically. The script's exact version is pinned in the Code Availability section; readers can re-execute it against the public audit log and reproduce every number in this paper.

%% =====================================================================
\section{The Challenge}
\label{sec:challenge}

\begin{table}[ht]
    \centering
    \small
    \setlength{\tabcolsep}{6pt}
    \begin{tabular}{lp{10cm}}
        \toprule
        \textbf{Term} & \textbf{Definition} \\
        \midrule
        \textbf{5-tuple credential}     & Structured equivalence claim $(r, T, \varepsilon, \alpha, \Pi)$: signal ratio, task, equivalence margin, significance level, patient subpopulation~\citep{ws2framework}. \\
        \textbf{Signal-equivalence}     & A reconstruction method is signal-equivalent to a named reference if its expected task performance on reduced-signal scans is within $\varepsilon$ of the reference's performance on full-signal scans, with confidence $\geq 1-\alpha$~\citep{ws2framework}. \\
        \textbf{Content-addressed hash} & A cryptographic fingerprint (e.g., SHA-256) of a dataset, framework, or method. Two artifacts with the same hash are guaranteed bit-identical. \\
        \textbf{RunBundle}              & PWM-format submission: pinned Docker image + eval entry point + published \texttt{results.json} + 5-tuple credentials. The scoring service re-executes the RunBundle and verifies the published results match the live execution. \\
        \textbf{Leave-one-vendor-out (LOVO)} & Cross-vendor evaluation protocol: train on data from $V - 1$ vendors, evaluate on the held-out vendor. \\
        \textbf{L2 / L3 / L4 spec}      & PWM registry layers: L2 = framework spec, L3 = dataset spec, L4 = method cert. Each layer is content-addressed and registered on the PWM blockchain. \\
        \textbf{Signal ratio} $r$       & Acquisition signal as fraction of reference: $r = $ mA-ratio (CT), $r = 1/$acceleration (MRI), $r = $ activity-ratio (PET). \\
        \textbf{Verdict}                & \texttt{PASS} / \texttt{FAIL} / \texttt{INDETERMINATE} on a 5-tuple credential, based on whether the bootstrap CI is contained in / disjoint from / straddles the equivalence band $(-\varepsilon, +\varepsilon)$. \\
        \bottomrule
    \end{tabular}
    \caption{Glossary of terms used throughout this review. Each is referenced multiple times; this table is provided as a single look-up for readers without a CS / informatics background.}
    \label{tab:glossary}
\end{table}

\subsection{Benchmark dataset (L3 spec)}

The challenge uses the PWM-LDCT dataset~\citep{ws1dataset}: \todo{N patient scans} paired-dose CT acquisitions across \todo{N vendors} vendor platforms at \todo{N sites} clinical sites. Each scan provides DICOM raw projections, image-space reconstructions at four dose levels (10\%, 25\%, 50\%, 100\%), and majority-vote radiologist annotations for lung-nodule detection and liver-lesion segmentation tasks. The dataset version evaluated in Year 1 is content-addressed at \texttt{pwm\_l3\_spec:pwm\_low\_dose\_ct\_benchmark:1.0.0}.

\subsection{Evaluation framework (L2 spec)}

Performance is reported through the signal-equivalence framework of~\citep{ws2framework}: every submission must declare its 5-tuple credentials of the form $(r, T, \varepsilon, \alpha, \Pi)$ for at least one canonical clinical task. The framework version evaluated in Year 1 is content-addressed at \texttt{pwm\_l2\_spec:dose\_equivalence\_v1:0.1.0}. Pixel-wise metrics (PSNR / SSIM / LPIPS) are reported as standard quality measures but are not the basis for leaderboard ranking; the leaderboard ranks by task-AUC on the canonical chest lung-nodule detection task.

\subsection{Submission protocol}

A submission is a PWM RunBundle: a content-addressed bundle containing a pinned Docker image, an eval entry point, a published \texttt{results.json}, and a corresponding 5-tuple credential JSON. The scoring service executes the RunBundle in a sandboxed container on the PWM-LDCT test split, verifies that the published \texttt{results.json} matches the live execution (within IEEE 754 FP tolerance), and posts the verified score. Failure to verify triggers an automatic rejection with a public failure-reason record. Tampered submissions are slashed via the standard PWM credential-violation protocol.

\subsection{Prize structure}

The Year-1 prize pool is \todo{N PWM tokens} (\$\todo{N USD equivalent}), distributed top-10:

\begin{table}[h]
    \centering
    \begin{tabular}{cc}
        \toprule
        \textbf{Place} & \textbf{Share} \\
        \midrule
        1st      & 25\% \\
        2nd      & 15\% \\
        3rd      & 10\% \\
        4th--10th & 50\% / 7 $\approx$ 7\% each \\
        \bottomrule
    \end{tabular}
\end{table}

We chose top-10 distribution over winner-take-all to encourage broad participation. Participation rates are reported in Section~\ref{sec:year1_cohort}.

%% =====================================================================
\section{Year-1 cohort}
\label{sec:year1_cohort}

\paragraph{Submission counts.} We received $N_{\textrm{sub}} = $ \todo{} community submissions between the challenge launch and the Year-1 close date. Of these, \todo{N} were from academic institutions, \todo{N} from industrial vendors, and \todo{N} from independent or unaffiliated submitters.

\paragraph{Methodology classes.} Submissions classified by architecture family in Table~\ref{tab:method_classes} \todo{populate}. Trend: \todo{e.g.\ diffusion-based reconstructions overtook transformer-based reconstructions in count}.

\begin{table}[h]
    \centering
    \begin{tabular}{lc}
        \toprule
        \textbf{Architecture family} & \textbf{Count} \\
        \midrule
        Unrolled iterative           & \todo{} \\
        Pure CNN (incl. encoder-decoder) & \todo{} \\
        Transformer-based            & \todo{} \\
        Diffusion / score-based      & \todo{} \\
        Hybrid / other               & \todo{} \\
        \midrule
        \textbf{Total}               & \todo{} \\
        \bottomrule
    \end{tabular}
    \caption{Year-1 submission distribution by architecture family.}
    \label{tab:method_classes}
\end{table}

\paragraph{Verification statistics.} Of the \todo{} submissions, \todo{} passed bit-identical re-execution verification on first attempt; \todo{} required resubmission to fix reproducibility issues; \todo{} were rejected for unrecoverable verification failure. The most common verification failure was \todo{e.g.\ missing CUDA RNG seed propagation}.

\paragraph{License posture.} \todo{N} submissions were released under Apache 2.0 or MIT licenses; \todo{N} under non-commercial or research-only licenses; \todo{N} did not declare a license (which is itself a leaderboard policy violation, flagged for the submitter).

%% =====================================================================
\section{Cross-method trends}
\label{sec:trends}

\subsection{Reconstruction quality}

PSNR / SSIM / LPIPS distributions are shown in Figure~\ref{fig:psnr_distribution}. Median PSNR at $r = 0.25$ across the cohort is \todo{X dB}; the top-3 submissions achieve \todo{Y dB}. The leaderboard top entry was \todo{submitter / method name} at \todo{Z dB}.

\todoFigure{fig:psnr_distribution}{Year-1 cohort PSNR / SSIM / LPIPS distributions on PWM-LDCT test split, per dose level}

\subsection{Downstream task performance}

Lung-nodule detection AUC at fixed FPR = 0.1 is summarized in Figure~\ref{fig:task_auc}. We observe that PSNR and task-AUC are \todo{strongly / weakly} correlated, with a Spearman $\rho = $ \todo{X} across the submission cohort. Notably, the top-PSNR submission is \todo{not / also} the top-task-AUC submission, illustrating that pixel-wise metrics are not interchangeable with clinical-task metrics.

\todoFigure{fig:task_auc}{Year-1 cohort lung-nodule detection AUC; scatter vs PSNR illustrating the pixel-wise vs task-metric gap}

\subsection{Cross-vendor generalization}

A subset of \todo{N} submissions reported cross-vendor results on the PWM-LDCT leave-one-vendor-out split. The median PSNR drop on the held-out vendor was \todo{X dB}; the best cross-vendor submission had a drop of only \todo{Y dB} (\todo{submitter / method}). The remaining \todo{N} submissions did not report cross-vendor results, which the leaderboard now records as a missing optional metric.

\subsection{Uncertainty quantification}

\todo{N submissions} reported per-pixel uncertainty; of these, \todo{N} provided calibration metrics (Spearman $\rho$ between predicted variance and observed error). The median calibration $\rho$ was \todo{X}. UQ remains underreported despite being one of the most clinically actionable supplementary outputs --- a recurring observation in Year 1 that we expect to inform Year-2 submission guidelines.

\subsection{Signal-equivalence credential coverage}

Per the framework of~\citep{ws2framework}, every submission was required to report at least one 5-tuple credential at $r = 0.25$. The distribution of credential verdicts is in Table~\ref{tab:credential_distribution}. The $r = 0.10$ regime --- a more aggressive dose reduction --- is currently unsolved by most submissions, with only \todo{N\%} of credentials passing at that level.

\todoTable{tab:credential_distribution}{Year-1 5-tuple credential verdict distribution (PASS / FAIL / INDETERMINATE) per dose ratio and per canonical task}

%% =====================================================================
\section{Clinical-impact analysis}
\label{sec:clinical_impact}

We summarize, for each canonical clinical task in the challenge, the most aggressive dose-reduction ratio at which at least one Year-1 submission achieves a passing signal-equivalence credential, broken down by patient subpopulation:

\begin{table}[h]
    \centering
    \footnotesize
    \begin{tabular}{lccc}
        \toprule
        \textbf{Task} & \textbf{Subpopulation} & \textbf{Best passing $r$} & \textbf{Method} \\
        \midrule
        Lung nodule detection $\geq 5$mm & Adult chest, lung screening   & \todo{}    & \todo{} \\
        Lung nodule detection $\geq 5$mm & Adult chest, oncology FU      & \todo{}    & \todo{} \\
        Liver lesion segmentation        & Adult abdomen                  & \todo{}    & \todo{} \\
        Liver lesion segmentation        & Adult oncology, post-chemo     & \todo{}    & \todo{} \\
        Body habitus quantification      & Pediatric chest (Year-2 stretch) & N/A     & --- \\
        \bottomrule
    \end{tabular}
    \caption{Most aggressive signal-equivalence-certified dose-reduction ratio by clinical task and subpopulation. ``Best passing $r$'' is the smallest $r$ at which at least one Year-1 submission reported a passing $(r, T, \varepsilon=0.02, \alpha=0.05, \Pi)$ credential.}
    \label{tab:clinical_impact_summary}
\end{table}

The clinical translation: a radiologist or physicist reading this table can map to the smallest dose at which their specific task (and their specific subpopulation) is currently supported by an open-source, cross-validated, signal-equivalence-credentialed reconstruction. Year-1 limitations are documented in Section~\ref{sec:open_problems}.

\subsection{Implications for clinical practice}

\paragraph{For radiologists.} The verdicts in Table~\ref{tab:clinical_impact_summary} are not regulatory clearance. They are evidence, in the same sense that a published reader study is evidence, that a specific open-source method, applied to the specified subpopulation and task, has been shown under bootstrap inference to be within $\varepsilon = 0.02$ AUC (or analogous metric) of the named full-dose reference, with 95\% confidence. Radiologists considering whether to lobby for protocol changes at their institution can cite specific 5-tuples rather than vendor marketing percentages.

\paragraph{For medical physicists.} The cross-vendor results in Section~\ref{sec:trends} provide a starting point for institution-specific QA. A method that achieves a passing credential when trained on the vendor a hospital uses, but a failing or indeterminate credential when trained on a different vendor and tested on this one, is a concrete starting point for the institution's commissioning plan.

\paragraph{For institutional review boards.} A protocol that proposes to reduce CT dose for a specified screening cohort can cite a passing 5-tuple credential as preliminary safety evidence for the diagnostic equivalence claim, replacing the present practice of citing vendor white papers with unspecified margins.

\paragraph{For regulators.} The challenge produces externally-verifiable, content-addressed evidence in a format aligned with the FDA's emerging guidance on AI/ML-based SaMD subgroup reporting~\citep{fda_aiml_2024}. A reviewer evaluating a submission can request the L4 cert hash and independently re-execute the published evaluation pipeline.

\subsection{Health equity and global health implications}

The PWM Low-Dose CT Challenge is structurally equity-positive along four axes that we believe are relevant to npj Digital Medicine's mandate around digital-health access.

\paragraph{Low-resource-institution access to validated methods.} Open-source weights, Apache 2.0 license, and a pip-installable evaluation library mean a hospital without an institutional license for vendor IR or a research-software budget can adopt a community-validated reconstruction method at zero marginal cost. Cross-vendor evaluation specifically addresses the scenario in which a low-resource institution operates a different scanner generation than the high-resource institutions where most methods are developed.

\paragraph{Reduced vendor lock-in.} A method that depends on vendor-proprietary IR pre-processing (closed-source kernels, vendor-specific noise modeling, vendor data formats) effectively forces an institution to keep buying from one manufacturer. The reference method and all qualifying leaderboard entries operate on raw sinograms and vendor-agnostic DICOM inputs, removing this structural lock-in.

\paragraph{Subgroup performance transparency.} Subgroup-stratified credentials make it harder to deploy a method that works well on a method-developer's training population but fails on under-represented subpopulations. A method missing a passing credential for $\Pi = \text{pediatric}$ or $\Pi = \text{specific-vendor-platform}$ cannot silently be marketed as "works for low-dose CT." This transparency is structurally pro-equity even before equity-specific subpopulations are added.

\paragraph{Global health perspective.} The Year-1 cohort includes \todo{N} clinical sites in \todo{N} countries; the Year-2 expansion roadmap (Section~\ref{sec:year2}) prioritizes international site recruitment. CT-based screening protocols (lung-cancer screening, tuberculosis screening) are increasingly deployed in middle-income countries; a credential framework that does not require vendor-specific or country-specific re-certification of every protocol lowers the per-site deployment cost.

We acknowledge that equity-positive structure is not the same as measured equity outcomes: whether the challenge actually narrows performance disparities depends on community participation, on which subpopulations are included in future dataset extensions, and on whether deployers act on the transparency the framework provides. We commit to reporting an explicit equity analysis in Year 2 once enough subpopulation-stratified data has accumulated to support it.

\subsection{Patient perspective}

From a patient's vantage point, the framework promises three concrete things, none of which depended on this challenge being framed as a digital-health intervention. (i) The same diagnostic confidence at lower radiation exposure for tasks where the credential PASSes at smaller $r$. (ii) Transparency about which method is recommended for the specific clinical context that applies to them: which vendor's scanner, which task, which subpopulation they belong to. A patient or patient-advocacy group can query the public registry by the parameters that match their own situation. (iii) An honest verdict of INDETERMINATE when the evidence does not support a definite answer, rather than a vendor's confident percentage backed by undisclosed evaluation conditions.

The current challenge does not yet include patient-reported outcomes (e.g., perceived anxiety from longer scan or repeat-scan protocols, willingness to consent to dose reduction); we list patient-reported outcomes as a Year-3 stretch addition once the Year-1 and Year-2 cohorts have established the technical baseline.

%% =====================================================================
\section{Open problems for Year 2}
\label{sec:open_problems}

We name three open problems we believe will define Year 2 of the challenge:

\paragraph{Open Problem 1: Pediatric subpopulations.} The Year-1 dataset is adult-only. Pediatric chest CT carries the highest per-scan radiation risk per unit dose and is where dose reduction matters most~\citep{larson2018pediatric}. The Year-2 dataset extension will add a pediatric chest subpopulation, and we call on the community to submit methods specifically trained or fine-tuned for it.

\paragraph{Open Problem 2: Cross-vendor at $r = 0.10$.} Methods that achieve passing credentials at $r = 0.25$ on the trained vendor often fail at $r = 0.10$ on a held-out vendor. We hypothesize this is a noise-statistics distribution-shift problem rather than an architecture problem; Year-2 submissions are encouraged to explicitly model the noise distribution across vendors.

\paragraph{Open Problem 3: Real-time inference for ER deployment.} The Year-1 top submissions average \todo{X seconds per slice} on an A100; emergency-radiology deployment requires \todo{Y seconds}. This is fundamentally an architecture-choice / quantization / distillation problem and is open for Year 2.

%% =====================================================================
\section{Year-2 challenge protocol}
\label{sec:year2}

The Year-2 challenge cycle opens on \todo{date} and closes on \todo{date}. Changes from Year 1:

\begin{enumerate}
    \item \textbf{Pediatric subpopulation added.} The dataset is extended to L3 spec v1.1.0; submissions report credentials for both adult and pediatric subpopulations.
    \item \textbf{Cross-vendor reporting becomes mandatory.} Every submission must report at least one cross-vendor leave-one-out result.
    \item \textbf{Inference-latency reporting is added to the leaderboard.} Latency is not a tiebreaker for the headline rank, but it is reported alongside performance metrics so users can pick by deployment constraint.
    \item \textbf{Industrial vendor invitation extended.} \todo{N vendors} were invited to Year 1; the Year-2 invitation list expands to all \todo{N vendors} major CT manufacturers.
\end{enumerate}

%% =====================================================================
\section{Discussion}

This paper inaugurates a publication mechanism that did not exist in the low-dose CT field before the launch of the PWM Challenge: a recurring, citable, peer-reviewed annual review whose content is reproducibly tied to a content-addressed benchmark and a content-addressed evaluation framework. Earlier surveys of low-dose CT reconstruction~\citep{wang2020deeplearningreview} have provided invaluable orientation to the field, but they are static snapshots; they cannot tell a reader whether a method published two years after the survey would still be ranked first on the same benchmark using the same evaluation protocol. The recurring annual review closes this gap.

We emphasize three structural properties of the challenge that we believe other medical-imaging fields should consider replicating.

\paragraph{Content-addressed benchmarks.} A benchmark version that can be silently re-extended, re-split, or re-cleaned between submissions makes year-over-year performance comparisons impossible. Anchoring the dataset version to a cryptographic hash registered on a public blockchain --- as the PWM Low-Dose CT Challenge does --- eliminates a class of benchmark-drift artifacts that have undermined longitudinal comparison.

\paragraph{Probabilistic, task-conditional evaluation.} A scalar PSNR or SSIM cannot capture the clinical question --- is the reconstruction good enough for the radiologist to make the diagnostic decision? --- and the gap between pixel-wise metrics and task performance is well-documented~\citep{kim2019smoothing}. The signal-equivalence framework of~\citep{ws2framework} provides a reporting standard that ties every claim to a specific task, subpopulation, margin, and confidence.

\paragraph{Re-executable submissions.} Bit-identical re-execution of every submission, verified by an automated scoring service, eliminates a separate class of reproducibility failures (cherry-picked random seeds, version-skew between training and reported eval, etc.). Adoption of this standard across medical-imaging benchmarks would substantially raise the field's reproducibility floor.

\paragraph{Comparison with benchmark-driven reviews in other medical fields.} Annual challenge reviews are not new to medical imaging --- the Brain Tumor Segmentation (BraTS) challenge has produced multi-year summaries~\citep{bakas2019brats}, the ISIC dermatology challenge tracks year-over-year skin-lesion methods~\citep{codella2019isic}, and the Kaggle RSNA pneumonia-detection challenge produced a sizable cited corpus. The PWM Low-Dose CT Challenge differs from these in three ways: (a) it uses a content-addressed benchmark hash anchored on a public registry rather than a centrally-administered dataset version, which closes a class of dataset-drift attacks that has affected prior challenge data~\citep{wang2020deeplearningreview}; (b) it requires a formal task-conditional credential per~\citep{ws2framework} rather than a single leaderboard scalar; (c) it commits in advance to a recurring annual review series rather than producing post-hoc summaries when an organizer has bandwidth. We expect each of these structural choices is independently transferable to non-CT challenges; we encourage benchmark organizers in other modalities and other clinical fields to consider adopting them.

\subsection*{Limitations of this review}

A review process is itself a form of evidence with its own biases. We acknowledge the following limitations of the present Year-1 synthesis.

\paragraph{Single-year cohort.} A Year-1 review cannot establish trends; longitudinal claims emerge only in Year-2 and later editions. We have framed the present review as a cohort summary plus open-problem identification rather than as a trend analysis precisely because we have no prior year against which to measure.

\paragraph{Organizer authorship.} This review is co-authored by the challenge organizers (see Conflict of interest in challenge organization, Section~\ref{sec:review_methods}). The safeguards adopted --- transparent ranking that does not put the reference method at the top, explicit author flagging in the leaderboard log, and a script-generated numeric content layer --- mitigate but do not eliminate the inherent risk that the organizers' perspective shapes the synthesis.

\paragraph{Clinical-task subset.} The two canonical tasks evaluated in Year 1 (lung-nodule detection, liver-lesion segmentation) cover an important but narrow slice of clinical CT use. Conclusions about ``low-dose CT'' as a whole drawn from this review should be qualified as ``conclusions about Year-1 challenge tasks''; thoracic-trauma, pulmonary-embolism, oncology-staging, cardiac-CT, and pediatric-CT readings are out of Year-1 scope.

\paragraph{Industrial-vendor participation.} Industrial-vendor submissions in Year 1 were lower than the academic-submitter rate (Table~\ref{tab:method_classes}). Conclusions about which method works best at the clinical-deployment frontier are partial until vendor participation increases; we explicitly invited the four major manufacturers and report participation outcomes.

\paragraph{Geographic representation.} The Year-1 dataset~\citep{ws1dataset} is collected at sites in \todo{N} countries, all high-income. International expansion is on the Year-2 roadmap; until then, performance claims should be qualified by population scope.

%% =====================================================================
\section*{Code and data availability}

The analysis script used to generate every numerical claim in this review is at \url{https://github.com/integritynoble/low_dose_CT/tree/main/WS-4_leaderboard/review/} (Apache 2.0). The script ingests the PWM Low-Dose CT Challenge audit log (publicly accessible at the leaderboard portal) and produces every table and figure deterministically. The exact pinned version used for this review is \todo{tagged-version-at-submission}. The underlying benchmark dataset is at~\citep{ws1dataset}, the signal-equivalence library at~\citep{ws2framework}, and the public leaderboard at the URL listed in Section~\ref{sec:challenge}.

%% =====================================================================
\section*{Acknowledgments}

We thank every Year-1 submitter for participating in the challenge and for their willingness to release their methods under permissive licenses. We thank the radiologists who contributed annotations to the PWM-LDCT dataset~\citep{ws1dataset}, the engineering team at the PWM Protocol Foundation for the scoring-service infrastructure, and the open-source maintainers of \texttt{pwm\_core}, \texttt{pwm\_ldct\_loader}, and \texttt{pwm\_dose\_equivalence}. Co-authors who are also authors of a submitted reconstruction method (initials: \todo{XX, YY}) recused themselves from the leaderboard-ranking computation per the COI protocol in Section~\ref{sec:review_methods}.

%% =====================================================================
\section*{Funding}

\todo{Declare at submission.} The PWM Low-Dose CT Challenge prize pool is funded by the PWM Protocol Foundation Reserve. Per-author funding sources (institutional research support, federal grants, etc.) are disclosed individually at submission per the \emph{npj Digital Medicine} funding-disclosure policy. The challenge does not accept industry sponsorship that is conditional on featuring or excluding any submission.

%% =====================================================================
\section*{Competing interests}

\todo{Refine at submission.} The authors declare the following structural conflicts:

\begin{itemize}
    \item Several authors are members of the team that organizes the PWM Low-Dose CT Challenge. The COI safeguards adopted by this review --- transparent ranking, recusal of author-of-submitted-method authors from leaderboard computation, organizer disclosure in the leaderboard log --- are documented in Section~\ref{sec:review_methods}.
    \item The PWM Protocol Foundation does not receive direct revenue from challenge participation or from any leaderboard entry. Foundation-affiliated authors hold no equity or royalty interest in any specific reconstruction method evaluated under the challenge.
    \item Individual authors hold the following per-person disclosures: \todo{list at submission}.
\end{itemize}

\section*{Author contributions}

\todo{Fill at submission per CRediT taxonomy: Conceptualization, Methodology, Software, Validation, Formal analysis, Investigation, Resources, Data curation, Writing-Original Draft, Writing-Review \& Editing, Visualization, Supervision, Project administration, Funding acquisition.}

%% =====================================================================
\section*{Reporting summary}

A \emph{Nature Portfolio} Reporting Summary in the format current at submission, covering subgroup-stratified performance reporting per~\citep{fda_aiml_2024}, code/data availability transparency, ethical-conduct of the challenge, and reproducibility of the review's analysis script, accompanies this manuscript as Supplementary Information.

%% =====================================================================
\bibliographystyle{plainnat}
\bibliography{refs}

\end{document}
